\documentclass{article}
\usepackage{mathtools} 
\usepackage{fontspec}
\usepackage[UTF8]{ctex}
\usepackage{amsthm}
\usepackage{mdframed}
\usepackage{xcolor}
\usepackage{amssymb}
\usepackage{amsmath}


% 定义新的带灰色背景的说明环境 zremark
\newmdtheoremenv[
  backgroundcolor=gray!10,
  % 边框与背景一致，边框线会消失
  linecolor=gray!10
]{zremark}{说明}

% 通用矩阵命令: \flexmatrix{矩阵名}{元素符号}{行数}{列数}
\newcommand{\flexmatrix}[4]{
  \[
  #1 = \begin{pmatrix}
    #2_{11}     & #2_{12}     & \cdots & #2_{1#4}   \\
    #2_{21}     & #2_{22}     & \cdots & #2_{2#4}   \\
    \vdots      & \vdots      & \ddots & \vdots     \\
    #2_{#31}    & #2_{#32}    & \cdots & #2_{#3#4}
  \end{pmatrix}
  \]
}

% 简化版命令（默认矩阵名为A，元素符号为a）: \quickmatrix{行数}{列数}
\newcommand{\quickmatrix}[2]{\flexmatrix{A}{a}{#1}{#2}}

\begin{document}
\title{习题 14.1}
\author{张志聪}
\maketitle

\section*{1}
\begin{itemize}
  \item (1)
        \begin{align*}
          \rho
           & = \lim\limits_{n\to\infty} \sqrt[n]{|a_n|} \\
           & = \lim\limits_{n\to\infty} \sqrt[n]{|n|}   \\
           & = 1
        \end{align*}
        于是收敛半径为$R = \frac{1}{\rho} = 1$。

        $x = 1$时，幂级数发散；$x = 1$时，幂级数收敛。
        故收敛区域为$(-1, 1)$。
  \item (2)
        \begin{align*}
          \rho
           & = \lim\limits_{n\to\infty} \sqrt[n]{|a_n|}                     \\
           & = \lim\limits_{n\to\infty} \sqrt[n]{|\frac{1}{n^2\cdot 2^n}|}  \\
           & = \lim\limits_{n\to\infty} \sqrt[n]{\frac{1}{n^2\cdot 2^n}}    \\
           & = \frac{1}{2} \lim\limits_{n\to\infty} \sqrt[n]{\frac{1}{n^2}} \\
           & = \frac{1}{2}
        \end{align*}
        于是收敛半径为$R = \frac{1}{\rho} = 2$。

        $x = 2$时，$\sum \frac{1}{n^2}$收敛；
        $x = -2$时，$\sum (-1)^n \frac{1}{n^2}$收敛。
        故收敛区域为$[-2, 2]$.

  \item (3)
        \begin{align*}
          \rho
           & = \lim\limits_{n\to\infty} \frac{|a_{n + 1}|}{|a_n|}                             \\
           & = \lim\limits_{n\to\infty} \frac{[(n + 1)!]^2}{[2(n + 1)]!} \frac{(2n)!}{(n!)^2} \\
           & = \lim\limits_{n\to\infty} \frac{(n + 1)^2}{(2n + 1)(2n + 2)}                    \\
           & = \frac{1}{4}
        \end{align*}
        于是收敛半径为$R = \frac{1}{\rho} = 4$。

        $x = 4$时，$\sum \frac{(n!)^2}{(2n)!}4^n$。
        \begin{align*}
          \lim\limits_{n\to\infty} \frac{(n!)^2}{(2n)!}4^n
           & = \lim\limits_{n\to\infty} \frac{(n!)^2 \cdot 2^n \cdot 2^n}{(2n)!}                                     \\
           & = \lim\limits_{n\to\infty} \frac{n! \cdot 2^n \cdot n! \cdot 2^n}{(2n)!}                                \\
           & = \lim\limits_{n\to\infty} \frac{(2n(2n - 2)\cdot 2) \cdot (2n(2n - 2)\cdot 2)}{2n(2n-1)(2n-2) \cdot 1} \\
           & = \lim\limits_{n\to\infty} \frac{2n(2n - 2)\cdot 2}{(2n-1)(2n-3) \cdot 1}                               \\
           & = \lim\limits_{n\to\infty} \frac{2n}{2n - 1} \frac{2n - 2}{2n - 3} \cdot \frac{2}{1}                    \\
           & \geq 1
        \end{align*}
        由级数收敛的柯西准则的推论可知，该级数发散；

        同理$x = -4$时，级数也发散。

        故收敛区域为$(-4, 4)$.

  \item (5)

        \begin{itemize}
          \item 方法一

                令
                \begin{equation*}
                  a_n =
                  \begin{cases}
                    \frac{1}{(2m - 1)!}, & n = 2m - 1 \\
                    0,                   & n = 2m
                  \end{cases}
                \end{equation*}
                于是原幂级数等价于
                \begin{align*}
                  \sum a_n (x - 2)^n
                \end{align*}
                我们有
                \begin{align*}
                  \rho
                   & = \overline{\lim\limits_{n \to \infty}} \sqrt[n]{|a_n|}         \\
                   & = \lim\limits_{m \to \infty} \sqrt[2m - 1]{\frac{1}{(2m - 1)!}} \\
                   & = \lim\limits_{t \to \infty} \sqrt[t]{\frac{1}{t!}}             \\
                   & = \lim\limits_{t \to \infty} \frac{1} {\sqrt[t]{t!}}            \\
                   & = 0
                \end{align*}
                故$R = +\infty$，收敛区域为$(-\infty, +\infty)$.

                \begin{zremark}
                  注
                  \begin{align*}
                    \lim\limits_{t \to \infty} \sqrt[t]{t!}
                     & = \lim\limits_{t \to \infty} e^{\ln \sqrt[t]{t!}}     \\
                     & = \lim\limits_{t \to \infty} e^{\ln (t!)^\frac{1}{t}} \\
                     & = \lim\limits_{t \to \infty} e^{\frac{1}{t} \ln t!}
                  \end{align*}
                  又
                  \begin{align*}
                    \lim\limits_{t \to \infty} \frac{\ln t!}{t}
                     & = \lim\limits_{t \to \infty} \frac{\ln 1 + \ln 2 + \cdots + \ln t}{t}
                  \end{align*}
                  由于$\{a_n\}, a_n = \ln n$
                  \begin{align*}
                    \lim \limits_{n \to \infty} a_n = +\infty
                  \end{align*}
                  由第二章总习题3可知
                  \begin{align*}
                    \lim\limits_{t \to \infty} \frac{\ln 1 + \ln 2 + \cdots + \ln t}{t} = +\infty
                  \end{align*}
                  综上
                  \begin{align*}
                    \lim\limits_{t \to \infty} \sqrt[t]{t!} = e^{+\infty} = +\infty
                  \end{align*}
                \end{zremark}

          \item 方法二：按照函数项级数。

                对任意$x$，有
                \begin{align*}
                  \lim\limits_{n\to\infty} \frac{|a_{n + 1}|}{|a_n|}
                   & = \lim\limits_{n\to\infty} \left|\frac{(x-2)^{2(n+1) - 1}}{(2(n+1)-1)!}\right| \left|\frac{(2n-1)!}{(x-2)^{2n - 1}}\right| \\
                   & = \lim\limits_{n\to\infty} \frac{(x - 2)^2}{2n(2n+1)}                                                                      \\
                   & = 0
                \end{align*}
                即对任意$x$级数都收敛，故$R = + \infty$，收敛区域为$(- \infty, \infty)$。

          \item 方法三：转换成不缺项幂级数级数。
                \begin{align*}
                  f_{2n - 1}'(x)
                   & = \frac{(x - 2)^{2n -2}}{(2n - 2)!}
                \end{align*}
                由定理14.7可知
                \begin{align*}
                  \sum \frac{(x - 2)^{2n - 1}}{(2n - 1)!} \, \text{和} \, \sum \frac{(x - 2)^{2n -2}}{(2n - 2)!}
                \end{align*}
                有相同的收敛半径。

                令$y = (x - 2)^2$，于是
                \begin{align*}
                  \sum \frac{(x - 2)^{2n -2}}{(2n - 2)!}
                   & = \sum \frac{y^{n - 1}}{(2n - 2)!}
                \end{align*}
                \begin{align*}
                  \rho
                   & = \lim\limits_{n \to \infty} \frac{|a_{n + 1}|}{|a_n|}                \\
                   & = \lim\limits_{n \to \infty} \frac{1}{(2(n+1) - 2)!} \cdot  (2n - 2)! \\
                   & = \lim\limits_{n \to \infty} \frac{1}{(2n - 1)2n}                     \\
                   & = 0
                \end{align*}
                故收敛半径为$R_y = +\infty$，即$y = (x - 2)^2 < +\infty$时，$\sum \frac{(x - 2)^{2n - 1}}{(2n - 1)!}$收敛，
                所以$\sum \frac{(x - 2)^{2n - 1}}{(2n - 1)!}$收敛半径为$R = +\infty$，
                收敛区域为$(- \infty, \infty)$.
        \end{itemize}

  \item (6)
        \begin{align*}
          \rho
           & = \lim\limits_{n \to \infty} \frac{|a_{n + 1}|}{|a_n|}                                               \\
           & = \lim\limits_{n \to \infty} \frac{3^{n+1} + (-2)^{n+1}}{n+1} \frac{n}{3^{n} + (-2)^{n}}             \\
           & = \lim\limits_{n \to \infty} \frac{3^{n+1} + (-2)^{n+1}}{3^{n} + (-2)^{n}} \frac{n}{n+1}             \\
           & = \lim\limits_{n \to \infty} \frac{3 + (-2)(-\frac{2}{3})^{n}}{1 + (-\frac{2}{3})^{n}} \frac{n}{n+1} \\
           & = 3
        \end{align*}
        所以$R = \frac{1}{\rho} = \frac{1}{3}$.
        \begin{align*}
          |x + 1| < \frac{1}{3} \\
          - \frac{4}{3} < x < - \frac{2}{3}
        \end{align*}
        $x = - \frac{2}{3}$时，
        \begin{align*}
          \sum \frac{3^n + (-2)^n}{n} (\frac{1}{3})^n
           & = \sum \frac{1 + (-\frac{2}{3})^n}{n}                \\
           & = \sum \frac{1}{n} + \sum \frac{(-\frac{2}{3})^n}{n}
        \end{align*}
        因为$\sum \frac{1}{n}$发散，$\sum \frac{(-\frac{2}{3})^n}{n}$收敛，
        故
        \begin{align*}
          \sum \frac{1}{n} + \sum \frac{(-\frac{2}{3})^n}{n}
        \end{align*}
        发散，即$x = - \frac{2}{3}$时，级数发散。

        $x = - \frac{4}{3}$时，
        \begin{align*}
          \sum \frac{3^n + (-2)^n}{n} (-\frac{1}{3})^n
           & = \sum \frac{(-1)^n + (\frac{2}{3})^n}{n}                \\
           & = \sum \frac{(-1)^n}{n} + \sum \frac{(\frac{2}{3})^n}{n}
        \end{align*}
        因为$\sum \frac{(-1)^n}{n}, \sum \frac{(\frac{2}{3})^n}{n}$都收敛，
        故$\sum \frac{(-1)^n}{n} + \sum \frac{(\frac{2}{3})^n}{n}$收敛，
        即$x = - \frac{4}{3}$时，级数收敛。

        综上，收敛区间为$[- \frac{4}{3}, - \frac{2}{3})$.

  \item (7)

        提示：用比式判别法。

  \item (8)
        \begin{align*}
          \lim\limits_{n \to \infty} \sqrt[n]{\left|\frac{x^{n^2}}{2^n}\right|}
           & = \lim\limits_{n \to \infty} \frac{|x|^n}{2}
        \end{align*}
        根据级数的根式判别法，当$\lim\limits_{n \to \infty} \frac{|x|^n}{2} < 1$时，级数收敛。

        $|x| = 1$时，$\lim\limits_{n \to \infty} \frac{|x|^n}{2} = \frac{1}{2}$；

        $|x| < 1$时，$\lim\limits_{n \to \infty} \frac{|x|^n}{2} = \frac{1}{2}$；

        $|x| > 1$时，$\lim\limits_{n \to \infty} \frac{|x|^n}{2} = +\infty$；

        故$R = 1$，收敛区间$[-1, 1]$。

\end{itemize}

\section*{2}

常见思路：
$n$在分母上，往往是对幂级数求导。

$n$不在分母上，从幂级数$1 + x + x^2 + \cdots + x^n + \cdots = \frac{1}{1 - x}$出发，
逐项求导，观察结果与目标幂级数的差异，完成证明。本题(3)是这个思路做的。
\begin{itemize}
  \item (1)

        对任意$x$，使用比式判别法：
        \begin{align*}
          \lim\limits_{n \to \infty} \frac{|u_{n + 1}(x)|}{|u_n(x)|}
           & = \lim\limits_{n \to \infty} \frac{x^{2(n+1) + 1}}{2(n+1) + 1} \cdot \frac{2n + 1}{x^{2n + 1}} \\
           & = \lim\limits_{n \to \infty} \frac{2n + 1}{2(n+1) + 1} \cdot \frac{x^{2(n+1) + 1}}{x^{2n + 1}} \\
           & = x^2
        \end{align*}
        故$x^2 < 1$，即$x \in (-1, 1)$上，级数绝对收敛。

        $x = 1$时，
        \begin{align*}
          1 + \frac{1}{3} + \frac{1}{5} + \cdots
          \leq \frac{1}{2} \sum \frac{1}{n}
        \end{align*}
        由比较原理可知
        $1 + \frac{1}{3} + \frac{1}{5} + \cdots$发散；
        $x = -1$时，也发散。

        故幂级数的收敛域为$(-1, 1)$.

        设幂级数的和函数为$f(x)$，即
        \begin{align*}
          f(x) = x + \frac{x^3}{3} + \frac{x^5}{5} + \cdots
        \end{align*}
        于是
        \begin{align*}
          f'(x) = 1 + x^2 + x^4 + \cdots + x^{2n} + \cdots
        \end{align*}
        以上是等比数列求和，故
        \begin{align*}
          f'(x) & = \lim\limits_{n \to \infty} \frac{1 - x^{2n + 1}}{1 - x^2} \\
                & = \frac{1}{1 - x^2}
        \end{align*}
        于是由逐项求积得到
        \begin{align*}
          f(x) & = \int_{0}^{x} \frac{1}{1 - t^2} dt                            \\
               & = \frac{1}{2}\int_{0}^{x} \frac{1}{1 + t} + \frac{1}{1 - t} dt \\
               & = \frac{1}{2} (\ln (1 + x) - \ln (1 - x))                      \\
               & = \frac{1}{2} \ln \frac{1 + x}{1 - x}
        \end{align*}

  \item (2)

        易得幂级数在$(-1, 1)$上一致收敛。

        设幂级数在$(-1, 1)$上一致收敛与$f(x)$，我们有
        \begin{align*}
          f(x) & = 1 \cdot 2 x + 2 \cdot 3 x^2 + \cdots + n(n + 1)x^n + \cdots       \\
          f(x) & = x(1 \cdot 2  + 2 \cdot 3 x + \cdots + n(n + 1)x^{n - 1} + \cdots)
        \end{align*}
        令
        \begin{align*}
          g(x) & = 1 \cdot 2  + 2 \cdot 3 x + \cdots + n(n + 1)x^{n - 1} + \cdots \\
          h(x) & = x^2 + x^3 + \cdots + x^n + \cdots
        \end{align*}
        于是
        \begin{align*}
          g(x) = h''(x)
        \end{align*}
        且
        \begin{align*}
          h(x)   & = \frac{x^2}{1 - x}   \\
          h''(x) & = \frac{2}{(1 - x)^3}
        \end{align*}
        所以
        \begin{align*}
          f(x) & = x g(x)                                     \\
               & = x h''(x)                                   \\
               & = \frac{2x}{(1 - x)^3} \, \, , x \in (-1, 1)
        \end{align*}

  \item (3)

        易得幂级数在$(-1, 1)$上一致收敛。

        设幂级数在$(-1, 1)$上一致收敛与$f(x)$。

        已知$1 + x + x^2 + \cdots + x^n + \cdots = \frac{1}{1 - x}$，
        逐项求导，我们有
        \begin{align*}
          g(x) = 1 + 2 x + 3 x^2 + \cdots + n x^{n - 1} + \cdots = (\frac{1}{1 - x})'
        \end{align*}
        又有
        \begin{align*}
          x g(x) = x + 2 x^2 + 3 x^3 + \cdots + n x^{n} + \cdots = x (\frac{1}{1 - x})'
        \end{align*}
        于是
        \begin{align*}
          (xg(x))' = 1 + 2^2 x + 3^2 x^2 + \cdots + n^2 x^{n - 1} + \cdots
        \end{align*}
        进而
        \begin{align*}
          x (xg(x))' = x + 2 x^2 + 3^2 x^3 + \cdots + n^2 x^{n} + \cdots
        \end{align*}
        得到目标函数，具体计算略。
\end{itemize}


\end{document}